% GENERATED FILE. Edit canonical lessons, then run npm run generate:books.
% canonical-source-hash: fnv1a64:5013eca5fe707567
% canonical-lessons: KA-C372-vanavasa, KA-C372-vidhisu, KA-C372-gadiparu, KA-C372-maji, KA-C372-kallatana

\chapter{Exile and Punishment}
\label{ch:ka-exile-and-punishment}
\hlchaptermodality{372}

\begin{chapteropening}
\textbf{By the end of this chapter:} \emph{I can say exile in the forest, to impose, banishment, former and theft.}

Complete the last lesson of chapter 372: I can say exile in the forest, to impose, banishment, former and theft.
\end{chapteropening}

\section[\texorpdfstring{\kn{ವನವಾಸ} (vanavāsa)}{ವನವಾಸ (vanavāsa)}]{\kn{ವನವಾಸ} (vanavāsa) — exile in the forest}

\label{lesson:KA-C372-vanavasa}



\begin{quote}
Before the new one: say the Kannada for a mansion, then the Kannada for to set foot.
\end{quote}

\subsection*{You'll want to know: \kn{ವನವಾಸ}}
\textbf{\kn{ವನವಾಸ}} — \emph{vanavāsa} — \textquotedblleft{}exile in the forest\textquotedblright{}.

\begin{quote}
\kn{ಕಥೆಯಲ್ಲಿ} \kn{ರಾಜಕುಮಾರ} \kn{ಹನ್ನೆರಡು} \kn{ವರ್ಷ} \kn{ವನವಾಸ} \kn{ಮಾಡಿದನು}.
\end{quote}

\begin{quote}
\emph{katheyalli rājakumāra hanneraḍu varṣa vanavāsa māḍidanu.}
\end{quote}

\begin{quote}
In the story the prince spent twelve years in exile in the forest.
\end{quote}

The epics are full of \kn{ವನವಾಸ}; today it can mean any long hard time away.

\begin{cousinweb}
\textbf{\kn{ವನವಾಸ}} (\emph{vanavāsa}) is borrowed from Sanskrit \emph{vanavāsa}, \textquotedblleft{}dwelling in the forest\textquotedblright{} (CDIAL 11265a), from \emph{vana} \textquotedblleft{}forest\textquotedblright{} and \emph{vāsa} \textquotedblleft{}dwelling\textquotedblright{}.
\end{cousinweb}

\begin{grammarlens}[title={What you've built}]
One more word for telling what happened.
\end{grammarlens}

\subsection*{Your turn}
\begin{itemize}
\raggedright
  \item \emph{Say it:} \emph{vanavāsa}
  \item \emph{Say it:} \emph{vanavāsa}, once more
  \item \emph{Say it:} \emph{katheyalli rājakumāra hanneraḍu varṣa vanavāsa māḍidanu.} — In the story the prince spent twelve years in exile in the forest.
  \item \emph{Recall:} say the Kannada for a mansion, then the Kannada for to set foot, then say \emph{vanavāsa} again
\end{itemize}

\begin{tcolorbox}[breakable,colback=teal!4,colframe=teal!35!black,title={Before you move on}]
What does \kn{ವನವಾಸ} mean? (\textquotedblleft{}Exile in the forest\textquotedblright{}.) Say it once more.
\end{tcolorbox}
\section[\texorpdfstring{\kn{ವಿಧಿಸು} (vidhisu)}{ವಿಧಿಸು (vidhisu)}]{\kn{ವಿಧಿಸು} (vidhisu) — to impose (a penalty)}

\label{lesson:KA-C372-vidhisu}



\begin{quote}
Before the new one: say the Kannada for to set foot, then the Kannada for exile in the forest.
\end{quote}

\subsection*{You'll want to know: \kn{ವಿಧಿಸು}}
\textbf{\kn{ವಿಧಿಸು}} — \emph{vidhisu} — \textquotedblleft{}to impose (a penalty)\textquotedblright{}.

\begin{quote}
\kn{ನ್ಯಾಯಾಧೀಶರು} \kn{ಕಳ್ಳನಿಗೆ} \kn{ದಂಡ} \kn{ವಿಧಿಸಿದರು}.
\end{quote}

\begin{quote}
\emph{nyāyādhīśaru kaḷḷanige daṇḍa vidhisidaru.}
\end{quote}

\begin{quote}
The judge imposed a fine on the thief.
\end{quote}

Past (I …): \kn{ವಿಧಿಸಿದೆ}, \emph{vidhiside}.

Punishments, fines, taxes and rules are imposed; the person takes -\kn{ಗೆ}.

\begin{cousinweb}
\textbf{\kn{ವಿಧಿಸು}} (\emph{vidhisu}) is built on Sanskrit \emph{vidhi}, \textquotedblleft{}a rule, an ordinance\textquotedblright{}, and \textquotedblleft{}fate\textquotedblright{} (CDIAL 11755), with the Kannada verb ending -\emph{isu}.
\end{cousinweb}

\begin{grammarlens}[title={What you've built}]
One more word for telling what happened.
\end{grammarlens}

\subsection*{Your turn}
\begin{itemize}
\raggedright
  \item \emph{Say it:} \emph{vidhisu}
  \item \emph{Say it:} \emph{vidhisu}, once more
  \item \emph{Say it:} \emph{nyāyādhīśaru kaḷḷanige daṇḍa vidhisidaru.} — The judge imposed a fine on the thief.
  \item \emph{Recall:} say the Kannada for to set foot, then the Kannada for exile in the forest, then say \emph{vidhisu} again
\end{itemize}

\begin{tcolorbox}[breakable,colback=teal!4,colframe=teal!35!black,title={Before you move on}]
What does \kn{ವಿಧಿಸು} mean? (\textquotedblleft{}To impose (a penalty)\textquotedblright{}.) Say it once more.
\end{tcolorbox}
\section[\texorpdfstring{\kn{ಗಡೀಪಾರು} (gaḍīpāru)}{ಗಡೀಪಾರು (gaḍīpāru)}]{\kn{ಗಡೀಪಾರು} (gaḍīpāru) — banishment; deportation}

\label{lesson:KA-C372-gadiparu}



\begin{quote}
Before the new one: say the Kannada for exile in the forest, then the Kannada for to impose.
\end{quote}

\subsection*{You'll want to know: \kn{ಗಡೀಪಾರು}}
\textbf{\kn{ಗಡೀಪಾರು}} — \emph{gaḍīpāru} — \textquotedblleft{}banishment; deportation\textquotedblright{}.

\begin{quote}
\kn{ನ್ಯಾಯಾಲಯ} \kn{ಆ} \kn{ರೌಡಿಗೆ} \kn{ಗಡೀಪಾರು} \kn{ಶಿಕ್ಷೆ} \kn{ವಿಧಿಸಿತು}.
\end{quote}

\begin{quote}
\emph{nyāyālaya ā rauḍige gaḍīpāru śikṣe vidhisitu.}
\end{quote}

\begin{quote}
The court sentenced that thug to banishment.
\end{quote}

\kn{ಗಡೀಪಾರು} \kn{ಮಾಡು}, to banish someone from a district or a country.

\begin{cousinweb}
\textbf{\kn{ಗಡೀಪಾರು}} (\emph{gaḍīpāru}) is usually traced to \kn{ಗಡಿ}, \textquotedblleft{}a border, a limit\textquotedblright{} (DEDR 1109), and \kn{ಪಾರು}, \textquotedblleft{}across, beyond\textquotedblright{}, from Sanskrit \emph{pāra}, \textquotedblleft{}the further bank\textquotedblright{} (CDIAL 8100): sent beyond the border.
\end{cousinweb}

\begin{grammarlens}[title={What you've built}]
One more word for telling what happened.
\end{grammarlens}

\subsection*{Your turn}
\begin{itemize}
\raggedright
  \item \emph{Say it:} \emph{gaḍīpāru}
  \item \emph{Say it:} \emph{gaḍīpāru}, once more
  \item \emph{Say it:} \emph{nyāyālaya ā rauḍige gaḍīpāru śikṣe vidhisitu.} — The court sentenced that thug to banishment.
  \item \emph{Recall:} say the Kannada for exile in the forest, then the Kannada for to impose, then say \emph{gaḍīpāru} again
\end{itemize}

\begin{tcolorbox}[breakable,colback=teal!4,colframe=teal!35!black,title={Before you move on}]
What does \kn{ಗಡೀಪಾರು} mean? (\textquotedblleft{}Banishment; deportation\textquotedblright{}.) Say it once more.
\end{tcolorbox}
\section[\texorpdfstring{\kn{ಮಾಜಿ} (māji)}{ಮಾಜಿ (māji)}]{\kn{ಮಾಜಿ} (māji) — former; ex-}

\label{lesson:KA-C372-maji}



\begin{quote}
Before the new one: say the Kannada for to impose, then the Kannada for banishment.
\end{quote}

\subsection*{You'll want to know: \kn{ಮಾಜಿ}}
\textbf{\kn{ಮಾಜಿ}} — \emph{māji} — "former; ex-".

\begin{quote}
\kn{ಮಾಜಿ} \kn{ಸೈನಿಕರು} \kn{ಸ್ಮಾರಕದ} \kn{ಬಳಿ} \kn{ಸೇರಿದರು}.
\end{quote}

\begin{quote}
\emph{māji sainikaru smārakada baḷi sēridaru.}
\end{quote}

\begin{quote}
Former soldiers gathered near the memorial.
\end{quote}

It stands before a noun: \kn{ಮಾಜಿ} \kn{ಶಿಕ್ಷಕ}, a former teacher.

\begin{cousinweb}
\textbf{\kn{ಮಾಜಿ}} (\emph{māji}) is a loan from Arabic \emph{māḍī}, \textquotedblleft{}past, gone by\textquotedblright{}, which came through Persian and Hindustani \emph{mājī}.
\end{cousinweb}

\begin{grammarlens}[title={What you've built}]
One more word for telling what happened.
\end{grammarlens}

\subsection*{Your turn}
\begin{itemize}
\raggedright
  \item \emph{Say it:} \emph{māji}
  \item \emph{Say it:} \emph{māji}, once more
  \item \emph{Say it:} \emph{māji sainikaru smārakada baḷi sēridaru.} — Former soldiers gathered near the memorial.
  \item \emph{Recall:} say the Kannada for to impose, then the Kannada for banishment, then say \emph{māji} again
\end{itemize}

\begin{tcolorbox}[breakable,colback=teal!4,colframe=teal!35!black,title={Before you move on}]
What does \kn{ಮಾಜಿ} mean? ("Former; ex-".) Say it once more.
\end{tcolorbox}
\section[\texorpdfstring{\kn{ಕಳ್ಳತನ} (kaḷḷatana)}{ಕಳ್ಳತನ (kaḷḷatana)}]{\kn{ಕಳ್ಳತನ} (kaḷḷatana) — theft}

\label{lesson:KA-C372-kallatana}



\begin{quote}
Before the new one: say the Kannada for banishment, then the Kannada for former.
\end{quote}

\subsection*{You'll want to know: \kn{ಕಳ್ಳತನ}}
\textbf{\kn{ಕಳ್ಳತನ}} — \emph{kaḷḷatana} — \textquotedblleft{}theft\textquotedblright{}.

\begin{quote}
\kn{ಊರಿನಲ್ಲಿ} \kn{ಕಳ್ಳತನ} \kn{ಹೆಚ್ಚಾದಾಗ} \kn{ಪೊಲೀಸರು} \kn{ಗಸ್ತು} \kn{ಶುರು} \kn{ಮಾಡಿದರು}.
\end{quote}

\begin{quote}
\emph{ūrinalli kaḷḷatana heccādāga polīsaru gastu śuru māḍidaru.}
\end{quote}

\begin{quote}
When thefts increased in the town, the police began patrols.
\end{quote}

\kn{ಕಳ್ಳ}, a thief, with -\kn{ತನ}, \textquotedblleft{}-ness, the doing of\textquotedblright{}: \kn{ಕಳ್ಳತನ} \kn{ಮಾಡು}, to steal.

\begin{cousinweb}
\textbf{\kn{ಕಳ್ಳತನ}} (\emph{kaḷḷatana}) is a native Dravidian noun: \emph{kaḷḷa}, \textquotedblleft{}a thief\textquotedblright{} (DEDR 1372), with the abstract ending -\emph{tana}.
\end{cousinweb}

\begin{grammarlens}[title={What you've built}]
One more word for telling what happened.
\end{grammarlens}

\subsection*{Your turn}
\begin{itemize}
\raggedright
  \item \emph{Say it:} \emph{kaḷḷatana}
  \item \emph{Say it:} \emph{kaḷḷatana}, once more
  \item \emph{Say it:} \emph{ūrinalli kaḷḷatana heccādāga polīsaru gastu śuru māḍidaru.} — When thefts increased in the town, the police began patrols.
  \item \emph{Recall:} say the Kannada for banishment, then the Kannada for former, then say \emph{kaḷḷatana} again
\end{itemize}

\begin{tcolorbox}[breakable,colback=teal!4,colframe=teal!35!black,title={Before you move on}]
What does \kn{ಕಳ್ಳತನ} mean? (\textquotedblleft{}Theft\textquotedblright{}.) Say it once more.
\end{tcolorbox}
